\section{CRM: oportunidades y embudo de ventas}

El \textbf{CRM} concentra el seguimiento comercial: cada trato con un prospecto o grupo se maneja como \textbf{oportunidad} y se visualiza en el \textbf{embudo} (\textbf{Pipeline} / vista Kanban por etapas).

\subsection{Abrir el CRM}

\begin{enumerate}
    \item En \textbf{Aplicaciones}, pulse \textbf{CRM}.
    \item Verá el menú lateral típico: \textbf{Oportunidades}, \textbf{Pipeline}, \textbf{Actividades}, \textbf{Informes}, etc. La vista principal suele ser el \textbf{Pipeline} (tarjetas por columnas = etapas).
\end{enumerate}

\subsection{Crear una oportunidad nueva}

\begin{enumerate}
    \item Pulse \textbf{Nuevo} (o \textbf{Crear} según la versión) en la vista de oportunidades, o use el \textbf{+} en una columna del embudo para crear la tarjeta ya en esa etapa.
    \item Complete el \textbf{Nombre} de la oportunidad (por ejemplo: ``Viaje Cusco -- familia Pérez -- abril'').
    \item En \textbf{Cliente}, busque y seleccione el \textbf{contacto} ya creado en el capítulo anterior, o cree uno rápido si el formulario lo permite con \textbf{Crear y editar}.
    \item Indique \textbf{Ingreso esperado} o montos si su flujo los usa en la oportunidad.
    \item Pulse \textbf{Guardar}.
\end{enumerate}

\subsection{Mover la oportunidad en el embudo}

\begin{enumerate}
    \item En la vista \textbf{Kanban} (tablero), arrastre la tarjeta de una columna a otra (por ejemplo: \textbf{Nuevo}, luego \textbf{Calificado}, \textbf{Propuesta}, \textbf{Ganado}), según las \textbf{etapas} que tenga su instancia.
    \item También puede abrir la oportunidad y cambiar la \textbf{Etapa} desde el formulario si su pantalla lo muestra.
\end{enumerate}

\subsection{Desde la oportunidad: cotización y seguimiento}

\begin{enumerate}
    \item Abra la oportunidad con un clic en la tarjeta o en el nombre en vista lista.
    \item Use el chatter para \textbf{Registrar una nota} o \textbf{Enviar mensaje} al equipo.
    \item Si aparece \textbf{Nueva cotización} o botón equivalente (\textbf{Cotización}, \textbf{Nuevo presupuesto}), úselo para generar el documento de ventas vinculado al cliente (véase el capítulo de Ventas).
\end{enumerate}

\subsection{Lista de oportunidades}

En el menú \textbf{Oportunidades} puede cambiar a vista \textbf{Lista} o \textbf{Gráfico} con los iconos junto al buscador. Use \textbf{Filtros} y \textbf{Agrupar por} para encontrar tratos por etapa, vendedor o fecha.
